Finite-context latent world models can store old evidence internally yet fail
to expose it to a later decision.  We study this as an interface problem rather
than an architecture label.  A valid memory claim must establish old-cue
demand, post-context retention, frozen-host exposure, and downstream use under
matched controls.  We instantiate this ladder with masked-evidence sidecars on
frozen DINO-WM and LeWM hosts, where PushT visual binding passes a
checkpointed downstream-use gate and DINO-WM PointMaze converts retained goal
state into fixed-controller execution gains.  Wall exposes the key failure
mode: carrier-prior memory remains high while host-output exposure collapses
with cue age.  We then remove the manual cue-feature handoff with a
self-supervised salient patch-set JEPA objective: target frames and patches
are mined without cue labels, cue crops, or key-frame annotations, and a
compact slot state predicts fixed generic color/texture latents.  On OGBench
PointMaze-large and Cube-single this non-manual objective passes all 30
registered cells over ages 4/8/15 and five seeds: full-memory balanced
accuracy is 0.92--1.00 while reset and recent-only controls stay near chance.
The result is a scoped method-and-audit contribution: persistent memory in
latent world models is a property of the target--carrier--host--consumer
interface, not a universal property of any single recurrent, state-space, or
fixed-trust module.
